% \iffalse meta-comment
%
% hit-thesis-bib.dtx 是参考文献页
%
% \fi
%
% \section{参考文献页}
%
% ^^A ======================================================================
% ^^A Module thesis-bib: 参考文献加载与 natbib/gbt7714 设置（hit-thesis）
% ^^A ======================================================================
%
% \changes{v3.2a}{2026/08/09}{引用样式与著录列表参数抽到 \file{hit-bibstyle.dtx}
%   与报告共用，这里只留学位论文自己的部分}
% 底下走 \pkg{gbt7714}。\cs{bibstyle@inline} 由它带，不用自己定义。加载完之后
% 引用样式与著录列表的参数在共享层 \file{hit-bibstyle.dtx}，排在本模块之前。
%
% \changes{v3.2a}{2026/08/16}{\pkg{gbt7714} 从类加载期挪到导言区末尾装。它内部
%   无条件拉 \pkg{natbib}，在类里装死了用户在导言区换文献方案的路}
% \pkg{gbt7714} 不在类加载期装，挪到导言区末尾（\texttt{begindocument/before}
% 钩子）。它内部无条件 \cs{RequirePackage}\marg{natbib}，而 \pkg{natbib} 与
% \pkg{biblatex} 一族互斥——在类里就装上，等于替用户把文献方案定死了。挪到
% 导言区末尾，用户导言区写的东西才有先说话的机会，这也是 thuthesis 的做法。
%
% \changes{v3.2a}{2026/08/16}{\cs{bibstyle@inline} 改用转接宏重定义，行内
%   引用的括号与分隔跟 \option{brackets-in-cite}/\option{comma-in-cite} 走，
%   形状照抄 \pkg{gbt7714} 自带的那份}
% 配置分两截：装载那一句挂 \texttt{begindocument/before}；其余全挂
% \texttt{package/gbt7714/after}。后者不能并进前者：\pkg{gbt7714} 会重定义
% \env{thebibliography}，本模块的重定义必须排在它之后才能赢，而且用户要是自己
% 在导言区提前装了 \pkg{gbt7714}，配置也得跟着那次装载走。
% \changes{v3.2a}{2026/08/16}{装载前先探 \pkg{biblatex}：用户设了
%   \option{bib}\texttt{!=\{backend!=biber\}} 或自己装了 \pkg{biblatex}，
%   \pkg{gbt7714} 就不再进来}
% 装载带一道侦测：\pkg{biblatex} 在（\option{backend}\texttt{=biber} 由类装的，
% 或用户自己装的，都一样）就不装 \pkg{gbt7714}——两族互斥，见
% \file{hit-bibstyle.dtx} 的守卫。
%    \begin{macrocode}
%<*thesis-bib>
\AddToHook { begindocument / before }
  {
    \@ifpackageloaded { biblatex }
      { }
      { \RequirePackage[sort&compress]{gbt7714} }
  }
\AddToHook { package / gbt7714 / after }
  {
\cs_gset:Npn \bibstyle@inline
  {
    \bibpunct { \hit@bib@cite@open: } { \hit@bib@cite@close: }
      { \hit@bib@cite@sep: } {n} {,} {,}
    \cs_gset:Npn \NAT@spacechar { \ }
  }
\citestyle{numerical}
%    \end{macrocode}
%
% \changes{v3.1a}{2022/05/17}{按照 \hitszthesis{} 生成深圳校区的参考文献样式 \texttt{hitszthesis.bst}。\issue[97]}
% \begin{macro}{\@biblabel}
% 设置参考文献列表的序号中的括号为中文括号。
% \changes{v3.1b}{2022/06/06}{将 ［］前加上 \cs{mbox\{\}} 来获得正确的序号长度}
% 在［\#1］前加上 \cs{mbox\{\}} 来获得正确的序号长度，否则会因为 \hologo{XeLaTeX}下的标点挤压而计算错误，这里应该整个包裹起来，否则换行后缩进会有问题。
% \changes{v3.1d}{2025/03/03}{按着 \issue[211] 来修改，将 \cs{mbox\{\}} 包住整个 label}
% \changes{v3.1d}{2025/03/03}{按照 \issue[215] 将哈尔滨硕博毕业论文的参考文献序号也改为全角}
%    \begin{macrocode}
\bool_lazy_or:nnT
  { \bool_if_p:N \g__hit_shenzhen_bool }
  {
    \bool_lazy_and_p:nn
      { \bool_if_p:N \g__hit_harbin_bool }
      {
        \bool_lazy_or_p:nn
          { \bool_if_p:N \g__hit_master_bool } { \bool_if_p:N \g__hit_doctor_bool }
      }
  }
  { \cs_set:Npn \@biblabel #1 { \mbox { ［#1］ } } }
%    \end{macrocode}
% \end{macro}
%
% \begin{environment}{thebibliography}
% 修改默认的 \env{thebibliography} 环境来适应间距的要求。
%    \begin{macrocode}
\RenewDocumentEnvironment{thebibliography}{m}{%
%    \end{macrocode}
%
% \changes{v3.1d}{2025/03/03}{更改完全角符号后的半角点号可能出现在行首，规避一下}
% 这里我直接贴一个 issue 的链接：\href{https://github.com/CTeX-org/ctex-kit/issues/701}{CTeX-Kit \#701}
% \changes{v3.2a}{2026/08/07}{FullRight引入了太多bug, 取消之 (bst文件中加入了nolinebreak)}
%    \begin{macrocode}
  % \xeCJKDeclareCharClass{FullRight}{"002E}%
  \hit@clear@page\phantomsection%
  \hit@appendix@chapter*{\bibname}[\hit@bibliography@heading@en]
  \normalsize
%    \end{macrocode}
% \changes{v3.1b}{2022/06/06}{修改深圳校区的 \cs{labelsep} 为 0pt}
%    \begin{macrocode}
  \hit@if@option@TF{shenzhen}
    { \hit@bib@list:nn {#1} {0pt} }
    { \hit@bib@list:nn {#1} {0.5em} }
  \sloppy\frenchspacing
%    \end{macrocode}
% \changes{v2.0.07}{2019/03/02}{添加flushbottom到thebibliography环境中}
%    \begin{macrocode}
  \flushbottom
%    \end{macrocode}
% \changes{v2.0.03}{2018/10/08}{添加参考文献分割开关}
%    \begin{macrocode}
  \hit@if@option@TF{splitbibitem}%
  {%
  \clubpenalty0
  \@clubpenalty \clubpenalty
  \widowpenalty0%
  \interlinepenalty-50%
  }%
  {%
  \clubpenalty4000
  \@clubpenalty \clubpenalty
  \widowpenalty4000%
  \interlinepenalty4000%
  }%
\sfcode`\.\@m}
{\cs_set_nopar:Npn \@noitemerr
  {\@latex@warning{Empty `thebibliography' environment}}%
\endlist\frenchspacing}
%    \end{macrocode}
% \end{environment}
%    \begin{macrocode}
\cs_set:Npn \NAT@cite #1#2#3
      {\ifNAT@swa\NAT@@open\if*#2*\else#2\NAT@spacechar\fi
        #1\NAT@@close\if*#3*\else\textsuperscript{#3}\fi\else#1\fi\endgroup}
%    \end{macrocode}
%
% \changes{v3.2a}{2026/08/16}{传统一线补一个 \cs{printbibliography}，正文里这一句
%   两条路线都认}
% \pkg{biblatex} 一线排文献表写 \cs{printbibliography}。传统一线也给一个同名
% 命令（\cs{providecommand}，装了 \pkg{biblatex} 时轮不到它），指到总出口。
% \cs{supercite} 同理：\pkg{biblatex} 核心自带，传统一线补一个别名——这条线
% 的 \cs{cite} 本来就是上标，语义正好。\cs{citep} 与 \cs{citet} 不用补，
% \pkg{natbib} 与 gb 系 cbx 各自都带（cbx 特意做过与 natbib 的防冲突处理）
% \cs{hit@bib@print:}，于是正文里一句 \cs{printbibliography} 两条路线通用——
% 前提是文献库写在 \option{struct} 的 \option{bibliography} 里。
% thuthesis 的主文档要求用户按后端二选一注释切换，这里不用。
%    \begin{macrocode}
\providecommand \printbibliography { \hit@bib@print: }
\providecommand \supercite { \cite }
%    \end{macrocode}
%
% \changes{v3.2a}{2026/08/16}{\option{compress-min} 在传统一线改走官方
%   \texttt{mincompress}，装载完把现值补落一次}
% \option{compress-min} 的现值补落给 \pkg{gbt7714.sty} 的 \texttt{mincompress}
% ——键（或旧类选项 \option{citetwo}）多半在装载之前就设了。之后再设由键码
% 自己落，见 \file{hit-keylist.dtx}。
%    \begin{macrocode}
\hit@bib@cite@compress@apply:
  }
%    \end{macrocode}
%
% \begin{macro}{\bibauthor}
% \changes{v3.2a}{2026/08/07}{控制参考文献列表中作者是否全大写和最大显示数量}
% \cs{bibauthor}\marg{作者表} 由 \file{.bst} 调用，按 \option{bibmaxauthor} 截断，
% 按 \option{bibauthorupper} 决定要不要全大写。两个类各留一份而没有并进
% \file{hit-bibstyle.dtx}：单个作者的排法用 \cs{exp\_args:NNe} 在定义时就把
% \option{bibauthorupper} 的取值展开定死，得排在 \cs{ProcessKeyOptions} 之后，
% 而共享层排在类骨架之前。
%    \begin{macrocode}
\exp_args:NNe \cs_new:Nn \g_bib_author_cs:n {
    \bool_if:NTF \g_bib_authorupper_bool {
        \exp_not:N \text_uppercase:n {#1}
    } {
        {#1}
    }
}
\NewDocumentCommand{\bibauthor}{m}{
  \int_zero:N \l_tmpa_int
  \exp_args:Ne \tl_map_inline:nn {
    \exp_args:Ne \tl_tail:n {#1}
  } {
    \int_compare:nNnT
      \l_tmpa_int = \g_bib_maxauthor_int {
      \prg_break:n {\tl_head:N {#1}}
    }
    \int_incr:N \l_tmpa_int
    \g_bib_author_cs:n {##1}
  }
  \prg_break_point:
}
%    \end{macrocode}
% \end{macro}
%
% \begin{macro}{\citeup}
% \changes{v3.2a}{2026/08/07}{\cs{citeup}（上标引用宏）从 book-floats 迁入 book-bib，归位至引用语义所属模块}
% 上标引用宏 \cs{citeup}，语义属引用，不属浮动体。
%    \begin{macrocode}
\NewDocumentCommand \citeup { m } { \textsuperscript { \cite {#1} } }
%    \end{macrocode}
% \end{macro}
%
% \subsection{biblatex 一线}
%
% \changes{v3.2a}{2026/08/16}{学位论文类的 biblatex 排版配置：标题、著录列表、
%   全角序号、\cs{citeup}，与传统一线同一套版面}
% \option{backend}\texttt{=biber} 时的排版配置。机制在
% \file{hit-bibstyle.dtx}：\pkg{biblatex} 装载完先跑共用的选项换算，再跑这里的
% \cs{hit@bib@biblatex@setup:}。目标是与上面 \env{thebibliography} 那一套排出
% 同一个版面，参数逐项对照着搬。
%
% \begin{macro}{\hit@bib@biber@heading:,\hit@bib@biber@env@begin:}
% 标题与著录列表。标题三句照抄 \env{thebibliography} 的开头。列表参数对照
% \cs{hit@bib@list:nn}：最宽标号那格 \pkg{biblatex} 自己量好放在
% \cs{labelnumberwidth} 里，不用 \cs{settowidth}；\cs{usecounter} 一句不搬，
% 序号由 \pkg{biblatex} 排，不走 \texttt{enumiv}。间距与分页惩罚与传统一线
% 逐项相同。
%    \begin{macrocode}
\cs_new_protected:Npn \hit@bib@biber@heading:
  {
    \hit@clear@page \phantomsection
    \hit@appendix@chapter * { \bibname } [ \hit@bibliography@heading@en ]
    \normalsize
  }
\cs_new_protected:Npn \hit@bib@biber@env@begin:
  {
    \list
      {
        \printtext [ labelnumberwidth ]
          { \printfield { labelprefix } \printfield { labelnumber } }
      }
      {
        \str_if_eq:VnT \g_bib_labelalign_str { left }
          { \RenewDocumentCommand \makelabel { m } { ##1 \hfill } }
        \dim_set:Nn \labelwidth { \labelnumberwidth }
        \hit@if@option@TF { shenzhen }
          { \dim_set:Nn \labelsep { 0pt } }
          { \dim_set:Nn \labelsep { 0.5em } }
        \dim_set:Nn \itemindent { 0pt }
        \dim_set:Nn \leftmargin { \labelsep + \labelwidth }
        \dim_add:Nn \itemsep { -0.8em }
        \dim_set:Nn \parsep { \bibparsep }
      }
    \sloppy \frenchspacing
    \flushbottom
    \hit@if@option@TF { splitbibitem }
      {
        \clubpenalty 0
        \@clubpenalty \clubpenalty
        \widowpenalty 0
        \interlinepenalty -50
      }
      {
        \clubpenalty 4000
        \@clubpenalty \clubpenalty
        \widowpenalty 4000
        \interlinepenalty 4000
      }
    \sfcode `\. \@m
  }
%    \end{macrocode}
% \end{macro}
%
% \begin{macro}{\hit@bib@biblatex@setup:}
% 全角序号的条件与上面 \cs{@biblabel} 是同一份：深圳，或哈尔滨的硕博。落实
% 走 \pkg{biblatex-gb7714-2015} 的架子：它把序号标签的处理器挂进了每个条目的
% 选项链（\texttt{gbbiblabel} 选项，处理器叫 \cs{gbbiblabelopt@<值>}），排版时
% 逐条重放，导言区里直接改 \cs{mkgbnumlabel} 会被重放盖掉（实测）。所以注册
% 一个自家的处理器、把选项值设成它，重放替我们装。另一半校区装半角处理器，
% 免得 \option{bib-punct}\texttt{=full} 的方括号补丁把序号标签也带成全角——
% 标签的全半角归校区管，跟传统一线的 \cs{@biblabel} 同一份条件。
% gb 系样式才有这一套，认 \cs{mkgbnumlabel} 在不在。
%
% \cs{citeup} 在 \pkg{biblatex} 一线就是 \cs{cite}——gb7714-2015 的 \cs{cite}
% 本来就是上标带方括号，再套 \cs{textsuperscript} 会叠两层。
% 末尾把 \option{compress-min} 的现值补落到 \texttt{gbrefcompress}
% 计数器上——键可能在 \option{backend} 之前、\pkg{biblatex} 还没装时就设了。
%    \begin{macrocode}
\cs_new:cpn { gbbiblabelopt@hitfullwidth }
  { \renewcommand { \mkgbnumlabel } [1] { \mbox { ［##1］ } } }
\cs_new:cpn { gbbiblabelopt@hithalfwidth }
  { \renewcommand { \mkgbnumlabel } [1] { \mbox { [##1] } } }
\cs_new_protected:Npn \hit@bib@biblatex@setup:
  {
    \cs_if_exist:NT \mkgbnumlabel
      {
        \bool_lazy_or:nnTF
          { \bool_if_p:N \g__hit_shenzhen_bool }
          {
            \bool_lazy_and_p:nn
              { \bool_if_p:N \g__hit_harbin_bool }
              {
                \bool_lazy_or_p:nn
                  { \bool_if_p:N \g__hit_master_bool }
                  { \bool_if_p:N \g__hit_doctor_bool }
              }
          }
          { \ExecuteBibliographyOptions { gbbiblabel = hitfullwidth } }
          { \ExecuteBibliographyOptions { gbbiblabel = hithalfwidth } }
      }
    \defbibheading { bibliography } [ \bibname ] { \hit@bib@biber@heading: }
    \defbibenvironment { bibliography }
      { \hit@bib@biber@env@begin: } { \endlist \frenchspacing } { \item }
    \RenewDocumentCommand \citeup { m } { \cite {##1} }
    \hit@bib@cite@compress@apply:
  }
%</thesis-bib>
%    \end{macrocode}
% \end{macro}
%
% \Finale
